\documentclass[10pt,twocolumn]{article}

\usepackage[margin=0.75in]{geometry}
\usepackage{amsmath}
\usepackage{amssymb}
\usepackage{booktabs}
\usepackage{array}
\usepackage{tabularx}
\usepackage{times}
\usepackage{url}

\pagestyle{plain}
\providecommand{\keywords}[1]{\par\smallskip\noindent\textbf{Keywords: }#1\par}
\newcolumntype{Y}{>{\raggedright\arraybackslash}X}

\title{Inducing Executable Ontology Workflow Skills from Procedural Documents}

\author{Anonymous Authors\\Anonymous Institution}

\begin{document}
\maketitle

\begin{abstract}
Procedural documents describe how domain objects are inspected, transformed, and coordinated through multiple operations. Existing process-extraction methods recover activities and control-flow relations, but their nodes are usually text spans rather than executable operations grounded in an ontology. We formulate ontology workflow skill induction: given an Object Layer, an ontology Function Layer, and an SOP or user manual, construct a typed state-transition workflow whose nodes reference ontology functions, variables reference ontology objects and properties, conditions are ontology predicates, and edges are constrained by effect--precondition compatibility. We propose a structure-aware framework that jointly grounds procedural units to functions and objects, constructs control and data flow, infers state transitions and exception paths, and validates the result through ontology reasoning and execution tests. The output is an executable Workflow Skill rather than a textual process summary or generic BPMN graph.
\end{abstract}

\keywords{Workflow induction, SOP understanding, ontology grounding, procedural documents, executable skills}

\section{Introduction}
Standard operating procedures, maintenance manuals, user instructions, and service playbooks encode reusable operational knowledge. They describe which operations should be performed, which objects those operations affect, what state enables each step, how information flows between steps, and how the procedure changes under exceptional conditions.

Automatic process extraction has progressed from identifying activities and actors to constructing procedural graphs and BPMN-like models~\cite{bellan2022pet,neuberger2024universal,du2024paged}. Agent-oriented work executes SOPs through LLM planning and tool calls~\cite{kulkarni2025agents,nandi2025sopbench}. These approaches leave a critical gap. A process graph can be structurally plausible while referring to operations that do not exist, passing the wrong object between steps, or violating the state semantics of the underlying system.

This paper studies workflow construction inside an operational ontology. The inputs are an Object Layer containing object types, properties, relations, and state predicates; a Function Layer containing ontology-bound atomic functions; and a procedural document. The output is an Ontology Workflow Skill. Every function-call node references a Function Layer identifier. Every variable has an ontology type. Conditions and invariants are predicates over object properties and relations. Data-flow edges connect function outputs to compatible downstream inputs. Control-flow edges are checked against function preconditions and effects.

The ontology is therefore not used as prompt background. It acts as the workflow's type system, state model, and composition constraint. Removing it should measurably degrade object consistency, typed data flow, state-transition correctness, or executability. If it does not, the method is only a renamed procedural-graph extractor.

The contributions are:
\begin{itemize}
    \item We define a target representation for ontology-grounded Workflow Skills with function nodes, typed object variables, state predicates, data flow, invariants, and exception paths.
    \item We propose a joint document--ontology--function induction method that grounds procedural units to executable functions and ontology objects.
    \item We validate generated workflows through ontology constraints, effect--precondition reasoning, and execution tests.
    \item We define an evaluation protocol that explicitly ablates the Object Layer, Function Layer, state predicates, typed data flow, and execution validation.
\end{itemize}

\section{Related Work}
\subsection{Process Element and Procedural Graph Extraction}
PET annotates activities, actors, gateways, and flow relations in process texts~\cite{bellan2022pet}. Universal Prompting improves process-element extraction using carefully specified LLM prompts~\cite{neuberger2024universal}. PAGED introduces a large procedural graph extraction benchmark covering sequential actions, constraints, and gateways~\cite{du2024paged}. These datasets provide direct baselines for structural extraction, but they do not require nodes to bind to a finite ontology Function Layer or variables to remain consistent with ontology objects.

\subsection{Workflow and Process Model Generation}
Process-model generation maps text to BPMN-like representations~\cite{kourani2024processmodeling}. BPMN captures tasks, events, and gateways, but it does not by itself specify enterprise object identities, ontology properties, callable function contracts, or machine-checkable state effects required for execution.

\subsection{Semantic Workflow and SOP Execution}
Semantic workflow systems represent services, resources, provenance, and process state with ontologies~\cite{kaefer2018linkedworkflows,norouzi2025processodp}. Agent-S and SOP-Bench evaluate long-horizon SOP execution by agents~\cite{kulkarni2025agents,nandi2025sopbench}. These works motivate ontology-aware workflow execution, but our task is upstream: induce the executable workflow skill from procedural documents and an ontology.

\section{Method}
\subsection{Representation}
Given Object Layer $O$, Function Layer $F$, and document $D$, the system constructs:
\[
W=\langle G,V,X,P,E,I,\Pi\rangle,
\]
where $G$ is the control-flow graph, $V$ are typed object variables, $X$ are data-flow bindings, $P$ are state predicates, $E$ are function effects, $I$ are invariants and exception paths, and $\Pi$ is evidence. A function-call node must reference exactly one Function Layer identifier. Required inputs must be supplied by workflow inputs, prior function outputs, user interactions, or ontology relation traversal.

\subsection{Structure-Aware Parsing}
The parser preserves headings, enumerations, tables, warning blocks, cross-references, and page offsets. It identifies goals, prerequisites, operations, observations, expected results, conditions, exception clauses, and recovery statements. Document structure constrains candidate scope but does not alone determine workflow edges.

\subsection{Joint Function and Object Grounding}
For each procedural unit, the system retrieves candidate pairs $(f,o)\in F\times O$ using action semantics, target-object compatibility, input/output properties, preconditions, effects, and neighboring context. This avoids selecting a lexically similar function that operates on the wrong object type. Unsupported procedural units are labeled as unresolved capabilities rather than invented functions.

\subsection{Typed Data Flow and State Transitions}
The method tracks object variables across the complete document. Function outputs are bound to downstream inputs only when ontology types and relation paths are compatible. Document conditions are normalized into predicates over ontology objects. Effects are propagated along the graph so that a downstream precondition must be satisfied by the incoming state or by an explicit observation.

\subsection{Control, Exception, and Invariant Induction}
The model predicts sequence, branch, join, loop, and exception edges. Warnings and policies become invariants over ontology predicates. Exceptions are linked to the smallest supported scope and may select an alternative function, compensation path, user interaction, or escalation.

\subsection{Validation and Repair}
The validator checks function existence, variable types, required parameters, object identity preservation, effect--precondition compatibility, invariant preservation, branch reachability, and repeated side effects. Validation failures trigger localized repair of the implicated node, edge, predicate, or binding.

\section{Experiments}
\subsection{Datasets}
Evaluation is layered because no public dataset annotates the complete target. PET and Universal Prompting datasets support activity, actor, gateway, and relation evaluation~\cite{bellan2022pet,neuberger2024universal}. PAGED supports procedural graph and constraint evaluation~\cite{du2024paged}. MyFixit and MSPT support maintenance or procedural actions, tools, parts, materials, and arguments. SOP-Bench supports downstream execution stress tests~\cite{nandi2025sopbench}. We add a Full Ontology Workflow benchmark that annotates Function IDs, object variables, state predicates, typed data flow, exceptions, invariants, and executable tests.

\subsection{Experimental Isolation}
The main experiment uses a gold Object Layer and gold Function Layer to isolate workflow induction. Additional settings replace them with an automatically constructed Function Layer from Paper 1, a noisy Object Layer, or an incomplete Function Layer. This separation is necessary: Paper 2's main claim cannot depend on Paper 1 being correct.

\subsection{Baselines and Ablations}
Baselines include direct document-to-workflow JSON, Universal Prompting plus heuristic graph construction, PAGED-style procedural graph extraction, document-to-BPMN generation, Function Library prompting without Object Layer, Object Layer prompting without Function Layer, full ontology prompting without validation, our method without state-transition reasoning, our method without typed data flow, our method without execution validation, and the full method.

\subsection{Metrics}
We report procedural-unit F1, node and labeled-edge F1, graph edit distance, gateway and exception F1, Function grounding accuracy, Object binding accuracy, ontology identifier validity, state-predicate F1, ontology violation rate, variable definition/use F1, parameter binding F1, type-correct binding rate, object identity consistency, required-input coverage, effect--precondition compatibility accuracy, invalid transition rate, goal-state reachability, executable workflow rate, hallucinated Function rate, forbidden side-effect rate, and downstream task success.

\subsection{Analysis}
The ontology contribution is supported only if removing ontology components degrades ontology correctness or executable utility. We partition outputs into graph-correct and ontology-correct, graph-correct but ontology-invalid, behaviorally equivalent but structurally different, and graph-invalid or non-executable. This prevents ordinary graph metrics from hiding semantic failures.

\section{Conclusion}
This paper makes the ontology operational inside workflow extraction. Function nodes, object variables, conditions, effects, data flow, invariants, and outputs are represented in ontology terms and validated as state transitions. The task is therefore ontology-grounded workflow synthesis, not generic procedural graph extraction.

\begin{thebibliography}{99}

\bibitem{bellan2022pet}
P. Bellan et al.
\newblock PET: An Annotated Dataset for Process Extraction from Natural Language Text Tasks.
\newblock In \emph{Business Process Management Workshops}, 2022.

\bibitem{neuberger2024universal}
J. Neuberger et al.
\newblock A Universal Prompting Strategy for Extracting Process Model Information from Natural Language Text using Large Language Models.
\newblock arXiv:2407.18540, 2024.

\bibitem{du2024paged}
Y. Du et al.
\newblock PAGED: A Benchmark for Procedural Graph Extraction from Documents.
\newblock arXiv:2408.03630, 2024.

\bibitem{kourani2024processmodeling}
H. Kourani et al.
\newblock Process Modeling With Large Language Models.
\newblock arXiv:2403.07541, 2024.

\bibitem{kaefer2018linkedworkflows}
T. Kaefer et al.
\newblock Linked Data Workflows.
\newblock arXiv:1804.05044, 2018.

\bibitem{norouzi2025processodp}
N. Norouzi et al.
\newblock Semantic Representation of Processes with Ontology Design Patterns.
\newblock arXiv:2509.23776, 2025.

\bibitem{kulkarni2025agents}
A. Kulkarni et al.
\newblock Agent-S: SOP-Guided Agent Automation.
\newblock arXiv:2503.15520, 2025.

\bibitem{nandi2025sopbench}
S. Nandi et al.
\newblock SOP-Bench: Benchmarking Agents on Standard Operating Procedures.
\newblock arXiv:2506.08119, 2025.

\end{thebibliography}

\end{document}
